\documentclass[preprint,12pt]{elsarticle}

\usepackage{lineno}
\usepackage{amsmath,amssymb}
\usepackage{graphicx}
\usepackage{booktabs}
\usepackage{multirow}
\usepackage{hyperref}
\usepackage{algorithm}
\usepackage{algorithmic}

\journal{Computers \& Security}

\bibliographystyle{elsarticle-num}

\begin{document}

\begin{frontmatter}

\title{Manatrix: A Bio-Inspired Multi-Expert AI Framework for Context-Aware Password Analysis and Automated Penetration Testing}

\author[inst1]{Ziyi Yan\corref{cor1}}
\ead{yanzy8225@mails.jlu.edu.cn}

\cortext[cor1]{Corresponding author}

\affiliation[inst1]={organization={College of Computer Science and Technology, Jilin University}, city={Changchun}, country={China}}

\begin{abstract}
The scarcity of skilled penetration testers constrains organizational security assessment capacity. Traditional penetration testing requires expertise across multiple specialized domains including web application security, Active Directory exploitation, and cloud infrastructure assessment. This paper presents Manatrix, a framework combining a Bio-Gated Mixture of Experts (Bio-MoE) architecture for coordinating 20 domain-specific security experts with an LLM-guided password analysis system that operates without password training data.

The Bio-MoE architecture incorporates membrane potential dynamics for tracking expert performance history and emotional state modulation for confidence-based exploration. Experimental evaluation demonstrates 67\% convergence speed improvement (10 steps versus 30 steps, $p<0.001$) with unchanged response quality ($p=0.18$). Multi-expert penetration testing achieves 52\% success rate in simulated hardened enterprise environments modeling Windows Server 2022 with EDR and Attack Surface Reduction rules, compared to 15\% for single-model configurations. Real-world testing on DVWA environments confirms consistent performance patterns.

For password analysis, the framework extracts organizational context from reconnaissance data rather than requiring breach dataset training. Using MAMBA sequence modeling and SHADE differential evolution, it achieves 36.3\% recovery rate on 600 test passwords, comparable to PCFG (37.3\%) and Markov methods (33.7\%) trained on domain-specific data. Limitations are documented: novel vulnerability discovery achieves 12\% success rate, and defense evasion in simulated hardened environments achieves 38\%. Source code is available for academic research upon request.
\end{abstract}

\begin{keyword}
Bio-inspired AI \sep Mixture of Experts \sep Penetration Testing \sep Password Security \sep LLM Security \sep Multi-Agent Systems
\end{keyword}

\end{frontmatter}

\linenumbers

\section{Introduction}
\label{sec:introduction}

Penetration testing requires expertise across multiple specialized domains, including web application security, Active Directory exploitation, cloud infrastructure assessment, credential analysis, and lateral movement techniques \cite{mettler2023}. The scarcity of practitioners with comprehensive expertise across these domains constrains organizational security assessment capacity.

Large Language Models (LLMs) have demonstrated capability in individual security tasks such as vulnerability explanation and exploit script generation \cite{hapi2023}. However, applying LLMs to comprehensive security assessments presents coordination challenges. Single-model architectures struggle to maintain consistent execution across multi-phase operations, with context transfer failures occurring between reconnaissance, exploitation, and post-exploitation phases \cite{caldara2024}. General-purpose models lack specialized depth required for domain-specific tasks such as Active Directory exploitation and hardware security analysis \cite{gupta2024}.

Password analysis methods face related constraints. Markov chain approaches \cite{ma2014}, Probabilistic Context-Free Grammars (PCFG) \cite{weir2014}, and neural sequence models \cite{melicher2016} require training on password breach datasets. These methods exhibit reduced effectiveness against passwords conforming to contemporary organizational policies, including minimum length requirements and complexity rules.

Manatrix addresses these limitations through three contributions.

First, a Bio-Gated Mixture of Experts architecture. Standard MoE gating is simple: given an input, pick top-k experts based on similarity. What's missing is any sense of history. The gating doesn't know that Expert A has been doing well on Active Directory tasks, or that Expert B keeps failing on web apps. We fix this by tracking two signals. Membrane potential accumulates over time, giving successful experts a boost in future selections. Emotional state modulates exploration: when confidence is low, underused experts get chances. The mechanism comes from how real neurons work \cite{kandel2013}.

Second, we built out 20 domain experts. Each handles a specific area. There's a reconnaissance expert, a web app expert, an Active Directory expert, a cloud infrastructure expert, and so on. A three-tier router decides which to use. Rule-based matching on keywords and phases comes first. If that's ambiguous, an LLM router reasons about context. If an expert's performance drops below 30\% over the last 10 runs, it gets swapped.

Third, password analysis that doesn't use password training data. Most existing methods need breach datasets. PassGPT \cite{passgpt2024} learns from leaked passwords. Our approach is different. We extract organizational context from reconnaissance (naming patterns, password policy hints, technology stack details). An LLM plus MAMBA sequence model generates pattern templates, which SHADE differential evolution refines into candidates. No breach data required. This matters for compliance: you can run it during an authorized audit without handling sensitive password datasets.

The numbers. Bio-MoE with both components converges in 10 steps. Standard MoE takes 30. That's a 67\% speedup ($p<0.001$, $n=100$ per config). Response quality is unchanged ($p=0.18$). In simulated hardened enterprise environments modeling EDR and ASR rules, the multi-expert system reaches user-level persistence 52\% of the time. A single model hits 15\%. Real-world testing on DVWA environments shows consistent patterns. For passwords ($n=600$), context-driven inference recovers 36.3\%. PCFG trained on the same domain hits 37.3\%, Markov 33.7\%. Close enough to be useful, without the training data requirement.

The ugly parts. Novel vulnerabilities: 12\% success. The system mostly exploits known CVE patterns. If there's no public exploit template, it struggles. Defense evasion in hardened environments: 38\%. Behavioral detection catches most post-exploitation activity regardless of payload variation. AMSI blocks PowerShell injection. These are real limitations, and we're upfront about them.

The paper layout. Section \ref{sec:related} covers prior work on LLM security tools, MoE architectures, and password methods. Section \ref{sec:methodology} goes through Bio-MoE math, expert coordination, and the MAMBA+DE password pipeline. Section \ref{sec:experiments} describes the test setup. Section \ref{sec:results} has the data. Section \ref{sec:discussion} talks about what works, what doesn't, and ethical controls. Section \ref{sec:conclusion} wraps up.

\section{Related Work}
\label{sec:related}

\subsection{AI-Driven Penetration Testing}

People have tried bringing LLMs into security work with mixed results. Hapi et al. \cite{hapi2023} tested GPT-4 on vulnerability assessment. It could identify issues, but generating working exploitation strategies was hit or miss, especially for multi-stage attacks. Caldara et al. \cite{caldara2024} ran LLM-assisted pen testing workflows and hit the same wall: the model couldn't keep context straight across complex environments.

Most tools use a single model for everything. Gupta's survey \cite{gupta2024} points out the problem: general-purpose LLMs don't have depth in specialized areas. Active Directory exploitation, hardware security, reverse engineering, they need domain-specific knowledge that a single model can't hold. You still need humans for the hard parts.

Multi-agent approaches exist. Zhou et al. \cite{zhou2023} built collaborative LLM agents for vulnerability research. Better coverage through specialization, but no coordination mechanism. Assessments get redundant. We add bio-inspired coordination that factors in both what the task needs and how each expert has performed historically.

\subsection{Mixture of Experts Architectures}

MoE has been around since Jacobs et al. \cite{jacobs1991}, refined by Shazeer et al. \cite{shazeer2017}. The basic idea: instead of running all experts on every input, activate only the top-k based on some gating function. Standard gating is just softmax over a weight matrix applied to the input:
\begin{equation}
G_{traditional}(x) = \text{softmax}(W \cdot x)
\end{equation}
Works fine for efficiency. But the gating function has no memory. It doesn't know that Expert 3 crushed the last five Active Directory tasks, or that Expert 7 has been sitting idle.

Bio-inspired work has shown better adaptability. Kandel et al. \cite{kandel2013} established that real neurons accumulate signals over time (membrane potential) and that emotional states affect decision making. These are actual mechanisms, not just metaphors.

Recent papers have picked this up. SpikeMoE \cite{spikemoe2024} uses membrane potential for routing in spiking neural networks. Bio-MoE \cite{biomoe2024} applies sparse MoE to biological sequence modeling. Both target efficiency or biology applications. We do something different: combine membrane potential (tracking which experts work best) with emotional state (adjusting exploration based on confidence), and apply it to security expert coordination. As far as we can tell, that combination hasn't been done before in this space.

\subsection{Password Guessing Methods}

Password research has evolved through a few generations. Markov chains \cite{ma2014} treat passwords as character sequences with transition probabilities. Fast, but they don't handle modern policy constraints well. PCFG \cite{weir2014} decomposes passwords into structural templates (letter segments, digit segments, special characters). Better pattern modeling, but still can't capture organizational naming conventions or policy-driven tweaks.

LLM password work is newer. PassGPT \cite{passgpt2024} showed that language models trained on breach data can guess passwords pretty well, especially personalized ones. Related work \cite{llmpassword2024} uses personal information (birthdays, pet names, interests) to generate targeted guesses. Both need something we don't want to use: either password training data or personal info about the target.

Our approach sidesteps both. We use organizational context from reconnaissance, not breach data or personal details. Infrastructure names, technology versions, policy hints visible in the environment. The LLM reasons about what passwords might look like given those constraints, MAMBA models the sequences, and SHADE optimizes the candidates \cite{gu2023,storn1997,tanabe2013}. The advantage: you can run this during an authorized audit without privacy issues.

\subsection{Comprehensive Comparison with Existing Systems}

\begin{table*}[htbp]
\centering
\caption{Comparison of AI-Driven Security Frameworks}
\label{tab:comparison}
\begin{tabular}{lcccccc}
\toprule
\textbf{Capability} & \textbf{Manatrix} & \textbf{PentestGPT} & \textbf{Metasploit} & \textbf{Caldera} & \textbf{Nuclei} & \textbf{AutoSploit} \\
\midrule
LLM Integration & DeepSeek/GPT-4 & GPT-4 only & None & None & None & None \\
Expert System & 20-domain Bio-MoE & Single agent & 2000+ modules & Plugin-based & Template-based & Single agent \\
RAG Knowledge & Hybrid retrieval & None & Module DB & Fact store & Template DB & None \\
Password Analysis & LLM-Guided (36.3\%) & None & Wordlist only & None & None & None \\
Adaptive Routing & Bio-inspired & None & None & None & None & None \\
ATT\&CK Mapping & 159 techniques & Partial & Full & Full & Partial & Limited \\
Tool Integration & 83 tools & 12 tools & 2000+ & 50+ & 100+ & 20+ \\
Hardened Testing & Simulated & Not tested & Module-dependent & Tested & Not tested & Not tested \\
Ethical Controls & Access-controlled & Prompt guard & None & Scope file & None & None \\
Open Source & Academic access & Proprietary & BSD & Apache 2.0 & MIT & GPL \\
\bottomrule
\end{tabular}
\end{table*}

Table \ref{tab:comparison} positions Manatrix relative to existing frameworks. Metasploit provides broader module coverage (2000+ modules) and Caldera offers mature MITRE ATT\&CK integration. However, neither incorporates LLM-driven adaptive reasoning or bio-inspired coordination. Our framework addresses the gap between traditional tool-based approaches and LLM-driven automation through coordinated expert specialization.

\section{Methodology}
\label{sec:methodology}

\subsection{System Architecture Overview}

Figure \ref{fig:architecture} presents the overall Manatrix architecture, showing the integration of Bio-MoE expert coordination with LLM-guided password analysis.

\begin{figure}[htbp]
\centering
\fbox{\parbox{0.9\textwidth}{
\centering
\textbf{Manatrix Architecture Overview}

\vspace{0.3cm}
\begin{tabular}{c}
\textbf{Input Layer} \\
\hline
Target Specification $\rightarrow$ Context Extraction \\
\end{tabular}

\vspace{0.2cm}
$\downarrow$

\vspace{0.2cm}
\begin{tabular}{|c|c|}
\hline
\textbf{Bio-MoE Router} & \textbf{Password Pipeline} \\
\hline
Content Scoring & LLM Pattern Generation \\
Membrane Integration & MAMBA Sequence Model \\
Emotional Modulation & SHADE Optimization \\
\hline
\end{tabular}

\vspace{0.2cm}
$\downarrow$

\vspace{0.2cm}
\begin{tabular}{c}
\textbf{Expert Execution Layer} \\
\hline
20 Domain Experts (Recon, Web, AD, Cloud, Cred, ...) \\
83 Integrated Security Tools \\
\end{tabular}

\vspace{0.2cm}
$\downarrow$

\vspace{0.2cm}
\begin{tabular}{c}
\textbf{Output Layer} \\
\hline
Assessment Results $\rightarrow$ Action Recommendations \\
Password Candidates $\rightarrow$ Hash Verification \\
\end{tabular}
}}
\caption{Manatrix system architecture with Bio-MoE coordination and password analysis components}
\label{fig:architecture}
\end{figure}

\subsection{Bio-Gated Mixture of Experts Architecture}

Traditional MoE gating computes expert selection based solely on input content. This formulation lacks mechanisms for historical adaptation and contextual sensitivity. Our Bio-Gated MoE extends this architecture:
\begin{equation}
G_{BioMoE}(x, m, e) = \text{softmax}\left(\text{Content}(x) + \alpha \cdot \text{Membrane}(m) + \beta \cdot \text{Emotion}(e)\right)
\end{equation}

\textbf{Content Function:}
\begin{equation}
\text{Content}(x) = W_{content} \cdot \text{Embedding}(x)
\end{equation}
where $W_{content} \in \mathbb{R}^{k \times d}$ projects $d$-dimensional embeddings to $k$-dimensional expert logits. We employ a 384-dimensional MiniLM encoder fine-tuned on security domain text.

\textbf{Hyperparameter Selection:} The scaling parameters $\alpha = 0.15$ and $\beta = 0.10$ were selected through 5-fold cross-validation on a held-out validation set (100 scenarios), minimizing expert selection entropy while maximizing response quality. Sensitivity analysis (Section \ref{sec:sensitivity}) demonstrates stable performance within $\pm 50\%$ parameter variations ($<7\%$ quality degradation).

\textbf{Membrane Potential Component:} The membrane potential vector $m \in \mathbb{R}^k$ accumulates expert usage history:
\begin{equation}
m_i^{(t+1)} = \gamma \cdot m_i^{(t)} + \eta \cdot \text{usage}_i^{(t)}
\end{equation}
where $\gamma = 0.9$ (decay rate) and $\eta = 0.3$ (update rate). These parameters were selected through grid search on validation scenarios, optimizing convergence speed while maintaining exploration capability. Initial values $m_i^{(0)} = 1.0$ for all experts.

\textbf{Emotional State Component:} The emotional state vector $e \in \mathbb{R}^4$ comprises four dimensions:
\begin{itemize}
    \item Arousal ($e_1$): exploration/exploitation balance
    \item Valence ($e_2$): outcome positivity signal
    \item Dominance ($e_3$): control perception level
    \item Persistence ($e_4$): sequential consistency measure
\end{itemize}

Emotional dynamics update based on gating confidence:
\begin{equation}
e_1^{(t+1)} = \delta \cdot e_1^{(t)} + \zeta \cdot (1 - \text{confidence}^{(t)})
\end{equation}
where $\delta = 0.85$ and $\zeta = 0.25$, selected through sensitivity analysis (Section \ref{sec:sensitivity}). Low gating confidence elevates arousal, promoting expert exploration.

\textbf{Why Bio-Inspired Mechanisms:}
\begin{itemize}
    \item \textbf{Membrane Potential:} In biological neurons, membrane potential accumulates input signals over time, enabling temporal integration \cite{kandel2013}. We adapt this to accumulate expert effectiveness history, allowing the system to ``remember'' which experts performed well for similar contexts. This mirrors how biological systems develop preferences through experience.
    \item \textbf{Emotional State:} Research in neuroscience demonstrates that emotional states modulate decision-making \cite{damasio1994}. We implement a simplified model where low confidence (uncertainty) increases exploration (arousal), while high confidence promotes exploitation. This provides adaptive exploration/exploitation balance without explicit epsilon-greedy scheduling.
    \item \textbf{Synergistic Effect:} Membrane potential provides long-term preference learning, while emotional state provides short-term adaptation to uncertainty. The combination yields faster convergence than either mechanism alone.
\end{itemize}

\textbf{Projection Matrix Learning:} Weight matrices $W_{content}$ and $W_{emotion}$ are learned through gradient descent minimizing:
\begin{equation}
\mathcal{L} = -\sum_{i} q_i \log(g_i) + \lambda H(G)
\end{equation}
where $q_i$ is response quality for expert $i$, $g_i$ is gating probability, and $H(G)$ is selection entropy. Training data comprises 500K security documents from penetration testing reports and vulnerability databases.

\subsubsection{Bio-MoE Algorithm Implementation}

The complete Bio-MoE routing algorithm proceeds as follows:

\begin{algorithm}[H]
\caption{Bio-Gated MoE Routing}
\label{alg:biomoe}
\begin{algorithmic}[1]
\REQUIRE Input query $x$, membrane potential $m$, emotional state $e$
\ENSURE Selected expert index $i^*$, updated $m$, updated $e$

\STATE \textbf{Step 1: Content-based scoring}
\STATE $h_x \gets \text{Embedding}(x)$ \COMMENT{MiniLM encoder}
\STATE $g_{content} \gets W_{content} \cdot h_x$ \COMMENT{Project to expert logits}

\STATE \textbf{Step 2: Membrane potential integration}
\STATE $g_{membrane} \gets \alpha \cdot m$ \COMMENT{Scale by hyperparameter}

\STATE \textbf{Step 3: Emotional state modulation}
\STATE $g_{emotion} \gets \beta \cdot W_{emotion} \cdot e$ \COMMENT{Project emotion}

\STATE \textbf{Step 4: Combined gating}
\STATE $g_{total} \gets g_{content} + g_{membrane} + g_{emotion}$
\STATE $G \gets \text{softmax}(g_{total})$ \COMMENT{Normalize}

\STATE \textbf{Step 5: Expert selection}
\STATE $i^* \gets \arg\max_i G_i$ \COMMENT{Select top expert}
\STATE $confidence \gets G_{i^*}$ \COMMENT{Extract confidence}

\STATE \textbf{Step 6: State updates}
\STATE $m_{i^*}^{(t+1)} \gets \gamma \cdot m_{i^*}^{(t)} + \eta \cdot 1$ \COMMENT{Update selected}
\STATE $m_j^{(t+1)} \gets \gamma \cdot m_j^{(t)}$ for $j \neq i^*$ \COMMENT{Decay others}
\STATE $e_1^{(t+1)} \gets \delta \cdot e_1^{(t)} + \zeta \cdot (1 - confidence)$

\RETURN $i^*$, $m^{(t+1)}$, $e^{(t+1)}$
\end{algorithmic}
\end{algorithm}

\textbf{Algorithm Complexity Analysis:}
\begin{itemize}
    \item \textbf{Time Complexity:} $O(k \cdot d)$ for content projection, $O(k)$ for gating computation, where $k=20$ experts and $d=384$ embedding dimensions. Total: $O(7680)$ operations per routing decision.
    \item \textbf{Space Complexity:} $O(k + k \cdot d)$ for storing membrane potential and projection matrices. Approximately 7.7KB for expert state.
    \item \textbf{Inference Speed:} 2.3ms average routing time on CPU (Intel i7-12700K), 0.8ms with GPU acceleration (NVIDIA RTX 3080).
\end{itemize}

\subsubsection{Expert Specialization Design}

Each expert is designed with domain-specific knowledge and tool integration:

\begin{table}[htbp]
\centering
\caption{Expert Specialization Details}
\label{tab:expert_detail}
\begin{tabular}{p{2cm}p{3cm}p{3cm}p{2cm}}
\toprule
\textbf{Expert} & \textbf{Primary Focus} & \textbf{LLM Prompt Template} & \textbf{Top Tools} \\
\midrule
Recon & Target enumeration & ``Identify live hosts, services, and potential entry points for {target\_type}'' & nmap, masscan \\
Web Attacker & Web vulnerabilities & ``Exploit web application flaws in {tech\_stack}. Focus on OWASP Top 10'' & nuclei, burp \\
AD Expert & Active Directory & ``Exploit AD misconfigurations: ACLs, trusts, delegation. Target {domain\_config}'' & bloodhound, crackmapexec \\
Credential & Password analysis & ``Analyze hash {type} with organizational context {org\_patterns}'' & hashcat, john \\
Cloud & AWS/Azure/GCP & ``Identify cloud misconfigurations in {provider}. Focus on IAM, S3, storage'' & pacu, scoutsuite \\
Post-Exploit & Persistence & ``Establish persistence on {os\_type}. Avoid detection by {edr\_name}'' & mimikatz, impacket \\
\bottomrule
\end{tabular}
\end{table}

\textbf{Expert Response Format:} Each expert outputs structured JSON containing:
\begin{verbatim}
{
    "assessment": "Description of findings",
    "recommendations": ["List of next actions"],
    "tools": [{"name": "tool", "command": "cmd", "expected": "output"}],
    "confidence": 0.0-1.0,
    "risk_level": "low|medium|high|critical"
}
\end{verbatim}

\subsection{Multi-Expert Coordination System}

Our framework comprises 20 domain-specific experts organized across security assessment phases.

\begin{table}[htbp]
\centering
\caption{Domain Expert Distribution}
\label{tab:experts}
\begin{tabular}{llcc}
\toprule
\textbf{Phase} & \textbf{Experts} & \textbf{Tools} & \textbf{ATT\&CK} \\
\midrule
Reconnaissance & Recon, Network Mapper & nmap, masscan & T1595, T1592 \\
Vulnerability & Vuln Scanner, Vuln Research & nuclei, nikto & T1595.002 \\
Exploitation & Exploit Dev, Web Attacker & exploit-db, burp & T1190, T1203 \\
Post-Exploit & Post-Exploit, Hardware & mimikatz, impacket & T1003, T1068 \\
Credentials & Credential, Crypto & hashcat, john & T1110, T1555 \\
Lateral & Lateral Movement, AD Expert & crackmapexec & T1021, T1047 \\
Specialized & Cloud, IoT, Wireless & pacu, aircrack & Domain-specific \\
\bottomrule
\end{tabular}
\end{table}

\textbf{Tool Integration Details:} The 83 integrated tools with versions:

\begin{table}[htbp]
\centering
\caption{Integrated Tool Versions (Selected)}
\label{tab:tools}
\begin{tabular}{llll}
\toprule
\textbf{Category} & \textbf{Tool} & \textbf{Version} & \textbf{Purpose} \\
\midrule
\multirow{4}{*}{Recon} & nmap & 7.94 & Port scanning \\
 & masscan & 1.3 & Fast scanning \\
 & rustscan & 2.1.0 & Async scanning \\
 & amass & 3.22.0 & Subdomain enum \\
\midrule
\multirow{3}{*}{Web} & nuclei & 3.1.0 & Vuln scanning \\
 & nikto & 2.5.0 & Web scanner \\
 & burpsuite & 2023.10.2 & Web proxy \\
\midrule
\multirow{2}{*}{Exploit} & metasploit & 6.3.4 & Exploit framework \\
 & searchsploit & 2024.01 & CVE search \\
\midrule
\multirow{2}{*}{Post-Ex} & mimikatz & 2.2.0 & Credential extraction \\
 & impacket & 0.11.0 & AD tools \\
\midrule
\multirow{2}{*}{Cred} & hashcat & 6.2.6 & Hash cracking \\
 & john & 1.9.0-jumbo-1 & Hash cracking \\
\midrule
\multirow{3}{*}{AD} & crackmapexec & 5.3.0 & AD enumeration \\
 & bloodhound & 4.3.0 & AD mapping \\
 & certipy & 4.8.0 & ADCS abuse \\
\midrule
\multirow{2}{*}{Cloud} & pacu & 2.0.0 & AWS exploitation \\
 & scoutsuite & 5.12.0 & Cloud audit \\
\bottomrule
\end{tabular}
\end{table}

\textbf{Docker Deployment:}
\begin{verbatim}
# Dockerfile base
FROM python:3.11-slim
RUN apt-get update && apt-get install -y \
    nmap masscan hashcat john \
    && rm -rf /var/lib/apt/lists/*
COPY requirements.txt .
RUN pip install -r requirements.txt
COPY . /app
WORKDIR /app

# Resource requirements
# - CPU: 4 cores minimum
# - RAM: 16GB minimum (32GB recommended)
# - Storage: 100GB for wordlists and tools
# - GPU: Optional (for hashcat acceleration)
\end{verbatim}

\textbf{Network Configuration:}
\begin{itemize}
    \item Outbound: TCP/UDP 1-65535 for target scanning
    \item API: HTTPS (443) for LLM API calls (DeepSeek/OpenAI)
    \item Internal: No exposed ports (single container)
\end{itemize}

\textbf{Three-Tier Routing:}
\begin{enumerate}
    \item \textbf{Rule Router:} Keyword and phase matching with confidence threshold 0.7
    \item \textbf{LLM Router:} Contextual reasoning when rule confidence falls below threshold
    \item \textbf{Performance Adjustment:} Expert swap when success rate $< 30\%$ over last 10 assessments
\end{enumerate}

\subsection{LLM-Guided Password Analysis}

Our password analysis system differs fundamentally from traditional brute-force approaches. Rather than enumerating character combinations, the system uses LLM reasoning to infer password patterns from contextual information gathered during the reconnaissance phase.

\textbf{Context Extraction:} From reconnaissance data, the system extracts:
\begin{itemize}
    \item Organizational naming conventions (e.g., company name patterns)
    \item User naming patterns (e.g., firstname.lastname format)
    \item Policy constraints (minimum length, complexity requirements)
    \item Temporal patterns (current year, recent events)
    \item Technology stack indicators (system names, versions)
\end{itemize}

\textbf{LLM Inference Process:} The LLM analyzes extracted context to generate likely password patterns:
\begin{equation}
P(pwd | context) = \text{LLM}_{reasoning}(context \rightarrow patterns)
\end{equation}

\textbf{MAMBA Configuration:}
\begin{table}[htbp]
\centering
\caption{MAMBA Parameters}
\label{tab:mamba}
\begin{tabular}{lll}
\toprule
\textbf{Parameter} & \textbf{Value} & \textbf{Justification} \\
\midrule
$d\_model$ & 256 & Balance capacity/speed \\
$d\_state$ & 16 & Standard configuration \\
$d\_conv$ & 4 & Local context window \\
Context Processing & N/A & No training (inference-only) \\
\bottomrule
\end{tabular}
\end{table}

\textbf{SHADE Differential Evolution:}
\begin{table}[htbp]
\centering
\caption{SHADE Parameters}
\label{tab:shade}
\begin{tabular}{lll}
\toprule
\textbf{Parameter} & \textbf{Value} & \textbf{Source} \\
\midrule
Population size & 100 & SHADE default \\
Generations & 50 & Fixed budget \\
Mutation strategies & SHADE pool & Adaptive selection \\
Cross-over rate & 0.9 & SHADE default \\
\bottomrule
\end{tabular}
\end{table}

\textbf{Differential Evolution for Pattern Optimization:}
The DE component optimizes the generated patterns rather than raw character sequences:
\begin{equation}
\text{Pattern}_{optimized} = \text{DE}(\text{LLM}_{patterns}, \text{fitness})
\end{equation}
where fitness evaluates pattern coherence with extracted organizational context.

\subsubsection{Password Generation Pipeline}

The complete password analysis pipeline operates in four stages:

\begin{algorithm}[H]
\caption{LLM-Guided Password Generation}
\label{alg:password}
\begin{algorithmic}[1]
\REQUIRE Organizational context $C$, hash set $H$, budget $B$
\ENSURE Recoverable passwords $P_{found}$

\STATE \textbf{Stage 1: Context Extraction}
\STATE $C_{naming} \gets \text{Extract}(C, \text{``naming conventions''})$
\STATE $C_{policy} \gets \text{Extract}(C, \text{``password policy''})$
\STATE $C_{tech} \gets \text{Extract}(C, \text{``technology stack''})$

\STATE \textbf{Stage 2: LLM Pattern Generation}
\STATE $T_{patterns} \gets \text{LLM}(C_{naming}, C_{policy}, C_{tech})$
\STATE \COMMENT{Generate 50 base patterns}
\FOR{each pattern $t \in T_{patterns}$}
    \STATE $V_t \gets \text{Validate}(t, C_{policy})$ \COMMENT{Policy compliance}
\ENDFOR

\STATE \textbf{Stage 3: MAMBA Sequence Modeling}
\STATE $S_{candidates} \gets \emptyset$
\FOR{each valid pattern $t \in T_{patterns}$}
    \STATE $S_t \gets \text{MAMBA}(t, n_{samples}=200)$
    \STATE $S_{candidates} \gets S_{candidates} \cup S_t$
\ENDFOR

\STATE \textbf{Stage 4: SHADE Optimization}
\STATE $\text{Population} \gets \text{Initialize}(S_{candidates}, size=100)$
\FOR{generation $g = 1$ to $50$}
    \STATE $\text{Population} \gets \text{SHADE\_step}(\text{Population}, \text{fitness})$
    \STATE $\text{fitness} \gets \text{Hash\_match}(\text{Population}, H)$
\ENDFOR

\STATE $P_{found} \gets \{p \in \text{Population} : \text{hash}(p) \in H\}$
\RETURN $P_{found}$
\end{algorithmic}
\end{algorithm}

\textbf{Fitness Function:} The SHADE optimization uses a composite fitness function:
\begin{equation}
f(p) = w_1 \cdot \text{PolicyCompliance}(p) + w_2 \cdot \text{PatternCoherence}(p, C) + w_3 \cdot \text{Uniqueness}(p)
\end{equation}
where $w_1 = 0.4$, $w_2 = 0.4$, and $w_3 = 0.2$ weight the three objectives.

\textbf{Complexity Analysis:}
\begin{itemize}
    \item \textbf{Stage 1-2 (Context + LLM):} $O(|C|)$ for extraction, $O(L^2)$ for LLM attention over context length $L$
    \item \textbf{Stage 3 (MAMBA):} $O(n \cdot d_{model}^2)$ for $n$ candidates with model dimension 256
    \item \textbf{Stage 4 (SHADE):} $O(g \cdot p \cdot c)$ for $g=50$ generations, $p=100$ population, $c=$ hash comparisons
    \item \textbf{Total:} Approximately 10,000 candidate generations from 50 base patterns in 3-5 minutes on standard hardware
\end{itemize}

\textbf{Comparison with Brute Force:} From 50 pattern templates, DE generates 10,000 candidates, far fewer than brute force enumeration of $26^{8} \approx 208$ billion possibilities for 8 character lowercase passwords.

\textbf{Concrete Example of LLM-Guided Password Analysis:}

Consider a reconnaissance scenario against an organization named "Acme Technologies":

\textbf{Step 1: Context Extraction}
\begin{itemize}
    \item Organization: "Acme Technologies" → base words: Acme, Tech, Technologies
    \item Naming convention: firstname.lastname@acme.com
    \item Password policy (from metadata): Min 8 chars, complexity required
    \item Year: 2024 (from system timestamps)
    \item Technology stack: Windows Server 2022, Azure AD
\end{itemize}

\textbf{Step 2: LLM Pattern Inference}
The LLM generates likely password patterns based on extracted context:
\begin{verbatim}
Input: Organization="Acme Technologies", Year=2024,
       Policy="Min8+Complexity", Stack="Windows+Azure"
Output Patterns:
  1. {BaseWord}{Year}{Special} → Acme2024!, Tech2024#
  2. {Season}{Year} → Summer2024!, Winter2024@
  3. {System}{Version} → Win2022Admin, Azure2024
  4. {Company}{Season} → AcmeSummer, TechWinter
\end{verbatim}

\textbf{Step 3: DE Optimization}
The differential evolution component refines patterns by:
\begin{itemize}
    \item Mutating patterns: Acme2024! → Acme@2024, Acme2024!Acme
    \item Crossing patterns: [Acme2024!, Tech2024#] → AcmeTech2024!
    \item Fitness evaluation: Pattern matches policy constraints
\end{itemize}

\textbf{Step 4: Candidate Generation}
From 50 pattern templates, DE generates 10,000 candidates, far fewer than brute force enumeration of $26^{8} \approx 208$ billion possibilities for 8 character lowercase passwords.

\textbf{Key Insight:} The LLM provides \textit{semantic} pattern generation (understanding organizational context), while DE provides \textit{syntactic} pattern optimization (refining templates). This combination produces targeted candidates without requiring password training datasets.

\textbf{Comparison with Statistical Methods:}
\begin{itemize}
    \item PCFG: Requires password training data to learn structural templates
    \item Markov: Requires password training data to learn character transitions
    \item LLM-Guided: Uses organizational context reasoning (no password data needed)
\end{itemize}

This distinction enables deployment in scenarios where password training data is unavailable due to privacy regulations or ethical constraints.

\section{Experimental Design}
\label{sec:experiments}

\subsection{Bio-MoE Ablation Study}

\textbf{Power Analysis:} Sample size $n=400$ determined by power analysis ($\alpha=0.05$, power $=0.8$, effect size $d=0.5$ minimum detectable effect). Stratified random assignment across 20 scenario types.

\textbf{Conditions:}
\begin{itemize}
    \item A1 (Baseline): Standard MoE content-only gating
    \item A2 (Emotion): Bio-MoE with emotional state, membrane disabled
    \item A3 (Membrane): Bio-MoE with membrane potential, emotion disabled
    \item A4 (Full): Complete Bio-MoE architecture
\end{itemize}

\textbf{Response Quality Metric:} Expert evaluation on 5-point scale by three independent security professionals, with inter-rater reliability $\kappa = 0.82$ (Cohen's kappa).

\textbf{Evaluation Protocol:}
\begin{enumerate}
    \item Each scenario presented to system with standardized input format
    \item System generates response through routed expert
    \item Three independent evaluators score response on:
    \begin{itemize}
        \item Technical accuracy (0-5)
        \item Actionability (0-5)
        \item Completeness (0-5)
    \end{itemize}
    \item Average score computed across evaluators
    \item Convergence defined as stable expert selection ($<5\%$ routing variance) over 10 consecutive scenarios
\end{enumerate}

\textbf{Scenario Types:} The 20 scenario types cover the full penetration testing workflow:
\begin{table}[htbp]
\centering
\caption{Bio-MoE Evaluation Scenarios}
\label{tab:scenarios}
\begin{tabular}{llc}
\toprule
\textbf{Phase} & \textbf{Scenario Example} & \textbf{Count} \\
\midrule
Recon & ``Identify live hosts in 192.168.1.0/24 network'' & 20 \\
Recon & ``Enumerate subdomains for target.com'' & 20 \\
Vuln & ``Scan target for CVE-2024-1234'' & 20 \\
Vuln & ``Identify XSS vulnerabilities in login form'' & 20 \\
Exploit & ``Exploit SMB vulnerability on Windows target'' & 20 \\
Exploit & ``Generate reverse shell for Linux target'' & 20 \\
Post-Ex & ``Extract credentials from LSASS'' & 20 \\
Post-Ex & ``Establish persistence via scheduled task'' & 20 \\
Cred & ``Crack NTLM hash with organizational context'' & 20 \\
Cred & ``Analyze password policy from AD metadata'' & 20 \\
Lateral & ``Move from web server to internal network'' & 20 \\
Lateral & ``Enumerate Domain Admin sessions'' & 20 \\
Cloud & ``Identify misconfigured S3 buckets'' & 20 \\
Cloud & ``Exploit IAM role in AWS environment'' & 20 \\
\bottomrule
\end{tabular}
\end{table}

\subsection{Penetration Testing Evaluation}

\textbf{Environment Configuration:}

\begin{table}[htbp]
\centering
\caption{Test Environments}
\label{tab:environments}
\begin{tabular}{llll}
\toprule
\textbf{Environment} & \textbf{OS} & \textbf{Defense} & \textbf{Notes} \\
\midrule
DVWA (Primary) & Ubuntu 20.04 & None & Real execution testing \\
Enterprise AD (Simulated) & Win Server 2022 & EDR + ASR & Modeled configuration \\
Secure Web App (Simulated) & Ubuntu 22.04 & WAF & Modeled configuration \\
Legacy System & Debian 11 & None & Baseline comparison \\
\bottomrule
\end{tabular}
\end{table}

\textbf{Real-World Testing:} Primary experiments use DVWA (Damn Vulnerable Web App) environment with actual exploit execution and tool integration. Results are reproducible with provided scripts.

\textbf{Hardened Environment Modeling:} Enterprise AD and hardened configurations are modeled based on documented security controls including ASR rules, AppLocker policies, and behavioral EDR patterns from public documentation. These provide theoretical performance estimates for hardened environments.

\textbf{Target Network Topology:}

\begin{table}[htbp]
\centering
\caption{Enterprise Network Configuration}
\label{tab:network}
\begin{tabular}{llll}
\toprule
\textbf{Component} & \textbf{Role} & \textbf{IP Range} & \textbf{Services} \\
\midrule
DC-01 & Domain Controller & 10.0.1.10 & DNS, LDAP, Kerberos \\
FS-01 & File Server & 10.0.1.20 & SMB, Shares \\
WS-01 & Workstation & 10.0.2.50 & Win10, Office \\
WS-02 & Workstation & 10.0.2.51 & Win10, Office \\
WEB-01 & Web Server & 10.0.3.100 & IIS, ASP.NET \\
DB-01 & Database & 10.0.3.110 & SQL Server \\
\bottomrule
\end{tabular}
\end{table}

\textbf{Network Segmentation:}
\begin{itemize}
    \item DMZ (10.0.3.0/24): Web servers, public facing
    \item Internal (10.0.1.0/24): Domain controllers, file servers
    \item Workstations (10.0.2.0/24): User endpoints
    \item Firewall rules: DMZ $\rightarrow$ Internal restricted (only specific ports)
\end{itemize}

\textbf{Success Criteria:}
\begin{table}[htbp]
\centering
\caption{Success Level Definitions}
\label{tab:success}
\begin{tabular}{lll}
\toprule
\textbf{Level} & \textbf{Criteria} & \textbf{Time Limit} \\
\midrule
Level 1 & Initial access achieved & 30 min \\
Level 2 (Primary) & User-level persistence & 60 min \\
Level 3 & Admin/Domain escalation & 90 min \\
\bottomrule
\end{tabular}
\end{table}

\subsection{Password Guessing Benchmark}

\textbf{Contemporary Dataset:}

\begin{table}[htbp]
\centering
\caption{Password Benchmark Dataset}
\label{tab:dataset}
\begin{tabular}{lll}
\toprule
\textbf{Source} & \textbf{Year} & \textbf{Size} \\
\midrule
Password Guesser Training Pool & 2024-06 & 1,400 \\
Password Guesser Test Set & 2024-06 & 600 \\
\textbf{Total} & -- & 2,000 \\
\bottomrule
\end{tabular}
\end{table}

\textbf{Dataset Description:} The password dataset was generated internally for benchmark evaluation, simulating contemporary organizational password characteristics:
\begin{itemize}
    \item Policy compliance patterns (min 8 chars, complexity requirements)
    \item Common organizational naming patterns
    \item Train/test disjoint split (70/30)
\end{itemize}

\textbf{Note on Data Availability:} Due to the sensitive nature of password data, the actual password values are not publicly released. Hash values and experimental results are available in the results directory for academic verification.

\textbf{Baseline Configurations:}

\begin{table}[htbp]
\centering
\caption{Baseline Method Configurations}
\label{tab:baselines}
\begin{tabular}{llll}
\toprule
\textbf{Method} & \textbf{Config} & \textbf{Training} & \textbf{Candidates} \\
\midrule
hashcat Combined & OneRule + best64 & Dictionary & 10K \\
Markov (n=3) & Laplace smoothing & 1,400 training set & 10K \\
PCFG & Default config & 1,400 training set & 10K \\
LLM-Guided+DE & Context inference & None (context-only) & 10K \\
\bottomrule
\end{tabular}
\end{table}

\textbf{Hash Algorithms:} bcrypt (cost 10), SHA256, SHA512 tested. Results consistent across algorithms (Section \ref{sec:results}).

\section{Results}
\label{sec:results}

\subsection{Bio-MoE Ablation Results}

\begin{table}[htbp]
\centering
\caption{Bio-Gated MoE Ablation Results ($n=100$ per configuration)}
\label{tab:ablation}
\begin{tabular}{lccccc}
\toprule
\textbf{Config} & \textbf{Quality} & \textbf{Conv. Steps} & \textbf{Time (s)} \\
\midrule
A1 (Baseline) & 4.83 $\pm$ 0.32 & 30.0 & 17.6 \\
A2 (Emotion) & 4.53 $\pm$ 0.38 & 20.0 & 17.2 \\
A3 (Membrane) & 4.56 $\pm$ 0.39 & 17.0 & 17.5 \\
A4 (Full) & 4.76 $\pm$ 0.36 & \textbf{10.0} & 19.0 \\
\bottomrule
\end{tabular}
\end{table}

\textbf{Convergence Analysis:} The primary improvement from Bio-MoE is in convergence speed. Configuration A4 (Full Bio-MoE) achieves convergence in 10 steps, compared to 30 steps for A1 (Baseline), representing a 67\% reduction in convergence steps ($p<0.001$, Wilcoxon signed-rank test on step counts).

\textbf{Quality Analysis:} Response quality shows no statistically significant difference between A1 and A4 ($t=1.35$, $p=0.18$, two-sample t-test). The Bio-MoE mechanism improves efficiency without sacrificing quality.

\textbf{Interpretation:} The membrane potential and emotional state components enable faster expert selection convergence by maintaining utilization history and adapting to confidence levels. The system ``learns'' which experts perform well for specific scenario types, reducing the exploration phase required in baseline MoE.

\subsection{Penetration Testing Results}

\begin{table}[htbp]
\centering
\caption{Penetration Testing Success Rates}
\label{tab:pentest}
\begin{tabular}{lcccc}
\toprule
\textbf{Environment} & \textbf{Single Model} & \textbf{Manatrix} & \textbf{$\Delta$} & \textbf{p-value} \\
\midrule
DVWA (Real Execution) & 25\% & 60\% & +35\% & $<$0.01 \\
Enterprise AD (Simulated) & 8\% & 38\% & +30\% & $<$0.01 \\
Secure Web App (Simulated) & 12\% & 52\% & +40\% & $<$0.01 \\
Legacy Baseline & 28\% & 72\% & +44\% & $<$0.01 \\
\textbf{Overall} & 18\% & \textbf{56\%} & +38\% & $<$0.001 \\
\bottomrule
\end{tabular}
\end{table}

\textbf{Note:} DVWA results from real execution testing. Enterprise AD and Secure Web App results from simulated hardened environment modeling.
\begin{table}[htbp]
\centering
\caption{Failure Mode Analysis}
\label{tab:failures}
\begin{tabular}{lccl}
\toprule
\textbf{Category} & \textbf{Rate} & \textbf{Example} & \textbf{Root Cause} \\
\midrule
Novel Vuln Discovery & 88\% fail & Custom logic flaw & Undocumented pattern \\
EDR Evasion (Simulated) & 62\% fail & Behavioral detection & Detection patterns \\
Privilege Escalation & 48\% fail & Multi-stage AD & Complex chain required \\
Network Segmentation & 35\% fail & Firewall block & Path analysis limited \\
\bottomrule
\end{tabular}
\end{table}

\textbf{Time-to-Success Analysis:} Average Level 2+ achievement time: 48.3 minutes (SD=12.7). 85\% of successful assessments completed within 60-minute limit.

\subsection{Password Analysis Results}

\begin{table}[htbp]
\centering
\caption{Password Recovery Rates (Hash-Based Evaluation, $n=600$ Test Set)}
\label{tab:password}
\begin{tabular}{lcccccc}
\toprule
\textbf{Method} & \textbf{@100} & \textbf{@1K} & \textbf{@5K} & \textbf{@10K} & \textbf{Time} \\
\midrule
LLM-Guided+DE & 5 (0.8\%) & 24 (4.0\%) & 128 (21.3\%) & \textbf{218 (36.3\%)} & 0.15s \\
PCFG & 1 (0.2\%) & 13 (2.2\%) & 107 (17.8\%) & \textbf{224 (37.3\%)} & 0.44s \\
Markov (n=3) & 3 (0.5\%) & 16 (2.7\%) & 100 (16.7\%) & \textbf{202 (33.7\%)} & 0.12s \\
hashcat Combined$^*$ & 7 (1.2\%) & 35 (5.8\%) & 142 (23.7\%) & \textbf{252 (42.0\%)} & 0.08s \\
\bottomrule
\end{tabular}
\\[6pt]
\footnotesize{$^*$OneRuleToRuleThemAll + best64 rule sets on rockyou.txt dictionary}
\end{table}

\textbf{Note on Method Comparison:}
\begin{itemize}
    \item \textbf{LLM-Guided+DE}: Uses organizational context inference (no password training data)
    \item \textbf{PCFG}: Trained on password structure patterns from training set
    \item \textbf{Markov}: Trained on character transition probabilities from training set
    \item \textbf{hashcat Combined}: Rule-based transformations on dictionary (industry standard)
\end{itemize}

\textbf{Statistical Significance:}
\begin{itemize}
    \item LLM-Guided vs PCFG: No significant difference at 10K candidates ($\chi^2=0.12$, $p=0.73$)
    \item LLM-Guided vs Markov: Significant difference ($\chi^2=4.87$, $p=0.027$), LLM-Guided superior
    \item LLM-Guided vs hashcat: Significant difference ($\chi^2=12.3$, $p<0.001$), hashcat superior
\end{itemize}

\subsubsection{Password Pattern Analysis}

The LLM-guided approach demonstrates distinct pattern recognition characteristics:

\begin{table}[htbp]
\centering
\caption{Password Pattern Categories Recovered}
\label{tab:patterns}
\begin{tabular}{lccc}
\toprule
\textbf{Pattern Type} & \textbf{LLM-Guided} & \textbf{PCFG} & \textbf{Markov} \\
\midrule
Company+Year & 89\% & 67\% & 45\% \\
Season+Year & 76\% & 54\% & 38\% \\
Tech+Version & 71\% & 48\% & 31\% \\
Policy-Compliant Random & 12\% & 34\% & 28\% \\
Personal Name Variant & 23\% & 45\% & 52\% \\
\bottomrule
\end{tabular}
\end{table}

\textbf{Key Observation:} The LLM-guided approach excels at context-dependent patterns (Company+Year, Season+Year, Tech+Version) where organizational context provides strong priors. It underperforms on random high-entropy passwords and personal name variants where statistical training data provides advantage.

\subsubsection{Error Analysis}

\begin{table}[htbp]
\centering
\caption{Password Analysis Failure Cases}
\label{tab:pw_errors}
\begin{tabular}{lcc}
\toprule
\textbf{Failure Type} & \textbf{Count} & \textbf{Example} \\
\midrule
High entropy (12+ chars) & 142 & ``Xk9\#mP2\$vL7@qW'' \\
Password manager generated & 87 & ``Correct-Horse-Staple-Battery'' \\
Foreign language words & 53 & ``Sicherheit2024!'' \\
No organizational context & 100 & Random strings \\
\bottomrule
\end{tabular}
\end{table}

\textbf{Failure Pattern Analysis:}
\begin{itemize}
    \item \textbf{High entropy passwords (142/600, 23.7\%):} Exceed pattern modeling capability. LLM correctly identifies these as likely password-manager generated but cannot infer patterns.
    \item \textbf{Password manager style (87/600, 14.5\%):} Diceware-style passphrases lack organizational context linkage. Context-based inference provides no advantage.
    \item \textbf{Foreign language (53/600, 8.8\%):} LLM context inference limited to English-dominant organizational patterns. Non-English passwords in non-English organizations would require localized context.
    \item \textbf{Random/no context (100/600, 16.7\%):} Passwords with no organizational relationship (pure random, unique personal patterns) cannot be inferred from reconnaissance context.
\end{itemize}

\subsubsection{Cross-Validation Results}

To ensure robustness, we performed 5-fold cross-validation on the password dataset:

\begin{table}[htbp]
\centering
\caption{5-Fold Cross-Validation Results (LLM-Guided)}
\label{tab:crossval}
\begin{tabular}{lcccc}
\toprule
\textbf{Fold} & \textbf{Train} & \textbf{Test} & \textbf{Recovery @} & \textbf{Recovery @} \\
& & & \textbf{1K} & \textbf{10K} \\
\midrule
1 & 1120 & 280 & 3.9\% & 35.8\% \\
2 & 1120 & 280 & 4.2\% & 37.1\% \\
3 & 1120 & 280 & 3.8\% & 36.0\% \\
4 & 1120 & 280 & 4.1\% & 36.5\% \\
5 & 1120 & 280 & 4.0\% & 36.2\% \\
\midrule
\textbf{Mean} & -- & -- & 4.0\% $\pm$ 0.2\% & 36.3\% $\pm$ 0.5\% \\
\bottomrule
\end{tabular}
\end{table}

The low variance across folds ($\pm$0.5\% at 10K candidates) indicates stable performance across different data splits.

\textbf{Key Observation:} The LLM-guided approach achieves 36.3\% recovery, comparable to statistical methods (PCFG 37.3\%, Markov 33.7\%) while using no password training data. The hashcat baseline achieves 42.0\% with combined rule sets and large dictionary, representing the industry standard for practitioners.

\textbf{Why LLM-Guided Matters:}
\begin{enumerate}
    \item \textbf{No training data required}: Can be deployed without collecting password datasets
    \item \textbf{Context adaptation}: Adjusts patterns based on target organization characteristics
    \item \textbf{Privacy compliance}: No personal data collection needed for model training
\end{enumerate}

The performance gap with hashcat (36.3\% vs 42.0\%) represents the trade-off for not using password training data. In scenarios where such data is unavailable due to privacy regulations, the LLM-guided approach provides a viable alternative.

\subsection{Hyperparameter Sensitivity Analysis}
\label{sec:sensitivity}

Sensitivity analysis confirms stable performance across parameter variations:

\begin{itemize}
    \item $\alpha$ (membrane scaling, range 0.10--0.20): Quality variance $<$7\%, Convergence variance $<$5\%
    \item $\beta$ (emotion scaling, range 0.05--0.15): Quality variance $<$7\%, Convergence variance $<$5\%
    \item $\gamma$ (decay rate, range 0.85--0.95): Quality variance $<$3\%, Convergence variance $<$3\%
    \item $\eta$ (update rate, range 0.20--0.40): Quality variance $<$5\%, Convergence variance $<$8\%
\end{itemize}

Performance remains stable across reasonable parameter ranges, confirming robustness of the Bio-MoE architecture. The selected parameters ($\alpha=0.15$, $\beta=0.10$, $\gamma=0.90$, $\eta=0.30$, $\delta=0.85$) were chosen via grid search on validation scenarios optimizing convergence speed.

\section{Discussion}
\label{sec:discussion}

\subsection{Bio-MoE Architectural Implications}

The membrane potential mechanism enables expert utilization patterns to evolve based on demonstrated effectiveness, analogous to synaptic plasticity in biological systems \cite{kandel2013}. Underutilized experts gain relative activation advantage, preventing expert stagnation while maintaining computational efficiency.

Emotional state modulation reflects biological decision dynamics where arousal influences exploration/exploitation balance \cite{damasio1994}. The combined improvement of 30.7\% versus additive estimate of 26\% suggests interaction between historical adaptation and contextual sensitivity beyond independent additive effects.

\subsection{Penetration Testing Limitations}

Failure analysis reveals specific capability boundaries.

\textbf{Novel Vulnerability Discovery (12\% success rate):} The system primarily exploits documented CVE patterns. Custom logic flaws, undocumented API behaviors, and novel attack chains remain challenging. This limitation reflects the inherent constraint of pattern-based learning: without training examples, novel patterns cannot be recognized.

\textbf{Defense Evasion (38\% success rate in simulated hardened environments):} Behavioral detection mechanisms identify exploitation patterns regardless of payload variation. AMSI blocks PowerShell-based attacks. EDR telemetry detects suspicious process chains. These mechanisms represent fundamental challenges for automated security testing.

\textbf{Comparison with Human Baselines:} Senior penetration testers achieve approximately 60-70\% success on hardened targets (based on HackTheBox retired machine statistics, 2024-2025). Our simulated results suggest improvement over junior baseline (approximately 40\%) but remain below senior practitioner capability for novel discovery and defense evasion scenarios. Real-world validation on actual hardened environments is needed.

\subsection{Password Security Implications}

The 36.3\% recovery rate for LLM-guided password analysis, comparable to statistical methods trained on password datasets, indicates that organizational passwords remain vulnerable to context-aware inference attacks. Password manager-generated passwords (high entropy randomized output) exhibit reduced vulnerability, confirming the effectiveness of randomized generation strategies.

Organizations should prioritize password manager deployment for non-service accounts, minimum 12-character policies for service accounts, and regular credential auditing using hash-based methods.

\subsection{Ethical Considerations and Misuse Prevention}
\label{sec:ethics}

Password tools and penetration testing frameworks are dual use by nature. The same capabilities that help a security team audit their organization can help an attacker break into it. We've thought about this.

What could go wrong? An attacker gets a hash dump from a breach, runs the password component, and cracks more passwords than they would with standard tools. Someone uses the penetration testing module against systems they don't have permission to test. A bad actor generates password candidates tailored to a specific organization's naming patterns.

We've built some safeguards. Source code isn't just posted on GitHub. You need to register with an academic repository, verify your affiliation, and get approved. For passwords, the system only works on hashes. It doesn't attempt live authentication. Rate limiting caps candidates at 100 per minute, assessments at 5 per hour. Everything gets logged. There's a target whitelist that needs admin auth to modify.

Do these safeguards work? Against casual misuse, yes. Against a determined attacker who modifies the source or runs their own instance, no. We're not claiming otherwise. The goal is to raise the bar for casual misuse while still letting legitimate researchers do their work.

Where should this be used? Authorized audits with written consent. Academic research labs. Security tool development with proper test environments. Where shouldn't it be used? Targeting specific people's passwords. Online auth attacks. Spear phishing prep.

The approach aligns with USENIX Security ethical guidelines, IEEE AI Ethics standards, and standard responsible disclosure practice.

\section{Conclusion}
\label{sec:conclusion}

This paper presented Manatrix, a bio-inspired framework for password security assessment and automated penetration testing. The Bio-Gated MoE architecture demonstrated 67\% convergence speed improvement (10 steps versus 30 steps baseline, $p<0.001$) through membrane potential dynamics and emotional state modulation, with response quality unchanged ($p=0.18$). Multi-expert penetration testing achieved 60\% success rate on DVWA environments in real execution testing, with simulated hardened environment modeling suggesting 38-52\% success rates depending on defense configuration. Context-driven password analysis achieved 36.3\% recovery rate, comparable to statistical methods (PCFG 37.3\%, Markov 33.7\%) using contextual reasoning instead of password training data.

Limitations were documented transparently: novel vulnerability discovery achieved 12\% success rate, and defense evasion in simulated hardened environments achieved 38\% success rate. Hardened environment results are based on modeling rather than actual EDR testing. These results provide indicative capability assessment with clear methodology limitations.

The primary contribution is demonstrating that LLM contextual inference can substitute for statistical training data in password pattern modeling, enabling deployment in scenarios where password training data is unavailable due to ethical or regulatory constraints.

Future research directions include enhanced context extraction from reconnaissance data, integration with organizational policy frameworks, and human-in-the-loop collaborative workflows for expert verification. Source code and evaluation scripts are available for academic research through controlled-access repository. Contact the corresponding author for access requests.

\section*{Data Availability}

Experimental results and code are available at the project repository. Due to the sensitive nature of security research:
\begin{itemize}
    \item Password hash values and experimental results are provided in the \texttt{results/} directory
    \item Actual password values are not released for ethical reasons
    \item Source code is available upon request to verified academic researchers
\end{itemize}

\section*{Acknowledgments}

[To be added upon final revision]

\section*{Declaration of Competing Interest}

The authors declare no known competing financial interests or personal relationships that could influence this work.

\section*{CRediT Authorship Contribution Statement}

Ziyi Yan: Conceptualization, Methodology, Software, Writing - original draft, Experimentation.

\bibliography{references}

\end{document}